\documentclass[10pt]{extarticle}
\usepackage[a4paper,margin=1.5cm]{geometry}
\usepackage[T1]{fontenc}
\usepackage{amsmath,amssymb}
\usepackage{graphicx}
\usepackage[export]{adjustbox}
\usepackage{multicol}
\usepackage{enumitem}
\usepackage{booktabs,array,tabularx}
\setlength{\columnsep}{0.7cm}
\setlength{\parindent}{0pt}
\setlength{\parskip}{3pt}
\pagestyle{empty}
\makeatletter
% figure registry (figdefs.tex is generated by the extract step)
\newcommand{\qfigdef}[4]{\@namedef{qf@file@#1}{#2}\@namedef{qf@w@#1}{#3}\@namedef{qf@wide@#1}{#4}}
\newcommand{\qf@check}[1]{\@ifundefined{qf@file@#1}{\PackageError{score2}{Unknown figure id #1}{}}{}}
\newcommand{\qf@img}[4][]{\includegraphics[width=#2\dimexpr\@nameuse{qf@w@#4}pt\relax,max width=#3,max height=0.45\textheight,#1]{\@nameuse{qf@file@#4}}}
% Full-width questions: finalize wraps every question that contains a wide figure in \qfullbegin...\qfullend,
% so the WHOLE question (stem, figures, options) is set across the page in reading order, then the two columns
% resume. The question's head - its number, stem and first figure (\qfullsplit marks the end of the head) -
% is kept on one page, so a number is never left at a page foot with its figure overleaf; the rest may break
% across pages normally, so no big gaps are left. Without a \qfullsplit, the whole question is kept together
% if it is under 45% of a page. Never write these yourself.
\newif\ifqf@full
\newif\ifqf@boxed
\newsavebox{\qf@box}
\newcommand{\qf@open}{\global\qf@fulltrue\par\global\qf@boxedtrue
  \setbox\qf@box\vbox\bgroup\hsize\textwidth\linewidth\textwidth}
\newcommand{\qf@emit}[1]{\par\egroup\global\qf@boxedfalse
  \ifdim\dimexpr\ht\qf@box+\dp\qf@box\relax<#1\textheight\noindent\usebox{\qf@box}\par\else\unvbox\qf@box\fi}
\newcommand{\qfullsplit}{\ifqf@boxed\qf@emit{0.85}\fi}
\newcommand{\qf@close}{\ifqf@boxed\qf@emit{0.45}\fi\global\qf@fullfalse}
\newcommand{\qfullbegin}{\end{multicols}\qf@open}
\newcommand{\qfullend}{\qf@close\begin{multicols}{2}}
\newcommand{\qfullstart}{\qf@open}
\newcommand{\qfulldone}{\qf@close}
% \qfig[scale]{id}: block figure, centred, directly where it is written (size set by the tool).
% Fallback if a wide figure is not inside a full-width question: break out of the columns in place - no floats.
\newcommand{\qfig}[2][1]{\qf@check{#2}%
  \ifqf@full\qf@blk{#1}{#2}\else
  \ifnum\@nameuse{qf@wide@#2}=1 \expandafter\@firstoftwo\else\expandafter\@secondoftwo\fi
  {\end{multicols}\begin{center}\qf@img{#1}{\textwidth}{#2}\end{center}\begin{multicols}{2}}%
  {\qf@blk{#1}{#2}}\fi}
\newcommand{\qf@blk}[2]{\par\smallskip{\centering\qf@img{#1}{\linewidth}{#2}\par}\smallskip}
% \qfigi[scale]{id}: inline figure (inside option lists / tables / minipages); never breaks columns.
\newcommand{\qfigi}[2][1]{\qf@check{#2}\qf@img[valign=c]{#1}{\linewidth}{#2}}
\makeatother
% structure macros - lint reads these to check completeness and figure placement
\newcommand{\papersec}[1]{\par\bigskip{\centering\Large\bfseries PAPER-#1\par}\medskip}
\newcommand{\partsec}[2]{\par\medskip\noindent{\large\bfseries PART-#1: #2}\par\smallskip}
\newcommand{\sect}[1]{\par\smallskip\noindent\textbf{#1}\par}
\newcommand{\qn}[1]{\par\smallskip\noindent\textbf{#1.}~}
\newcommand{\paraq}[2]{\par\smallskip\noindent\textbf{Paragraph for Questions #1 and #2}\par}
\newenvironment{qoptions}{\begin{itemize}[leftmargin=1.8em,itemsep=1pt,topsep=2pt,label={}]}{\end{itemize}}
\newcommand{\opt}[2]{\item[(#1)] #2}
% figure registry (generated by score2_tools.py finalize)
\qfigdef{p1_phy_s1_q01_1}{Score Advanced - Test 22 (10 Apr 2024)-img/p1_phy_s1_q01_1.png}{97.5}{0}
\qfigdef{p1_phy_s1_q03_1}{Score Advanced - Test 22 (10 Apr 2024)-img/p1_phy_s1_q03_1.png}{67.3}{0}
\qfigdef{p1_phy_s1_q04_1}{Score Advanced - Test 22 (10 Apr 2024)-img/p1_phy_s1_q04_1.png}{97.9}{0}
\qfigdef{p1_phy_s1_q04_A}{Score Advanced - Test 22 (10 Apr 2024)-img/p1_phy_s1_q04_A.png}{81.6}{0}
\qfigdef{p1_phy_s1_q04_B}{Score Advanced - Test 22 (10 Apr 2024)-img/p1_phy_s1_q04_B.png}{81.1}{0}
\qfigdef{p1_phy_s1_q04_C}{Score Advanced - Test 22 (10 Apr 2024)-img/p1_phy_s1_q04_C.png}{76.8}{0}
\qfigdef{p1_phy_s1_q04_D}{Score Advanced - Test 22 (10 Apr 2024)-img/p1_phy_s1_q04_D.png}{76.2}{0}
\qfigdef{p1_phy_s1_q05_1}{Score Advanced - Test 22 (10 Apr 2024)-img/p1_phy_s1_q05_1.png}{125.6}{0}
\qfigdef{p1_phy_s2_q07_1}{Score Advanced - Test 22 (10 Apr 2024)-img/p1_phy_s2_q07_1.png}{116.1}{0}
\qfigdef{p1_phy_s2_q09_1}{Score Advanced - Test 22 (10 Apr 2024)-img/p1_phy_s2_q09_1.png}{76.8}{0}
\qfigdef{p1_phy_s2_q11_1}{Score Advanced - Test 22 (10 Apr 2024)-img/p1_phy_s2_q11_1.png}{100.3}{0}
\qfigdef{p1_phy_s3_q03_1}{Score Advanced - Test 22 (10 Apr 2024)-img/p1_phy_s3_q03_1.png}{95.0}{0}
\qfigdef{p1_phy_s3_q03_2}{Score Advanced - Test 22 (10 Apr 2024)-img/p1_phy_s3_q03_2.png}{91.8}{0}
\qfigdef{p1_phy_s3_q04_1}{Score Advanced - Test 22 (10 Apr 2024)-img/p1_phy_s3_q04_1.png}{100.6}{0}
\qfigdef{p1_chem_s1_q03_1}{Score Advanced - Test 22 (10 Apr 2024)-img/p1_chem_s1_q03_1.png}{121.7}{0}
\qfigdef{p1_chem_s1_q03_A}{Score Advanced - Test 22 (10 Apr 2024)-img/p1_chem_s1_q03_A.png}{121.3}{0}
\qfigdef{p1_chem_s1_q03_B}{Score Advanced - Test 22 (10 Apr 2024)-img/p1_chem_s1_q03_B.png}{245.2}{0}
\qfigdef{p1_chem_s1_q03_C}{Score Advanced - Test 22 (10 Apr 2024)-img/p1_chem_s1_q03_C.png}{159.7}{0}
\qfigdef{p1_chem_s1_q03_D}{Score Advanced - Test 22 (10 Apr 2024)-img/p1_chem_s1_q03_D.png}{152.6}{0}
\qfigdef{p1_chem_s1_q04_1}{Score Advanced - Test 22 (10 Apr 2024)-img/p1_chem_s1_q04_1.png}{171.1}{0}
\qfigdef{p1_chem_s1_q04_A}{Score Advanced - Test 22 (10 Apr 2024)-img/p1_chem_s1_q04_A.png}{64.2}{0}
\qfigdef{p1_chem_s1_q04_B}{Score Advanced - Test 22 (10 Apr 2024)-img/p1_chem_s1_q04_B.png}{69.6}{0}
\qfigdef{p1_chem_s1_q04_C}{Score Advanced - Test 22 (10 Apr 2024)-img/p1_chem_s1_q04_C.png}{61.5}{0}
\qfigdef{p1_chem_s1_q04_D}{Score Advanced - Test 22 (10 Apr 2024)-img/p1_chem_s1_q04_D.png}{43.2}{0}
\qfigdef{p1_chem_s1_q05_C}{Score Advanced - Test 22 (10 Apr 2024)-img/p1_chem_s1_q05_C.png}{125.8}{0}
\qfigdef{p1_chem_s2_q10_1}{Score Advanced - Test 22 (10 Apr 2024)-img/p1_chem_s2_q10_1.png}{153.1}{0}
\qfigdef{p1_chem_s2_q10_2}{Score Advanced - Test 22 (10 Apr 2024)-img/p1_chem_s2_q10_2.png}{199.4}{0}
\qfigdef{p1_chem_s2_q12_1}{Score Advanced - Test 22 (10 Apr 2024)-img/p1_chem_s2_q12_1.png}{227.5}{0}
\qfigdef{p1_chem_s2_q13_1}{Score Advanced - Test 22 (10 Apr 2024)-img/p1_chem_s2_q13_1.png}{100.4}{0}
\qfigdef{p1_chem_s2_q13_D}{Score Advanced - Test 22 (10 Apr 2024)-img/p1_chem_s2_q13_D.png}{160.7}{0}
\qfigdef{p1_chem_s3_q05_1}{Score Advanced - Test 22 (10 Apr 2024)-img/p1_chem_s3_q05_1.png}{31.1}{0}
\qfigdef{p2_phy_s1_q01_1}{Score Advanced - Test 22 (10 Apr 2024)-img/p2_phy_s1_q01_1.png}{99.2}{0}
\qfigdef{p2_phy_s1_q01_A}{Score Advanced - Test 22 (10 Apr 2024)-img/p2_phy_s1_q01_A.png}{69.3}{0}
\qfigdef{p2_phy_s1_q01_B}{Score Advanced - Test 22 (10 Apr 2024)-img/p2_phy_s1_q01_B.png}{73.5}{0}
\qfigdef{p2_phy_s1_q01_C}{Score Advanced - Test 22 (10 Apr 2024)-img/p2_phy_s1_q01_C.png}{72.5}{0}
\qfigdef{p2_phy_s1_q01_D}{Score Advanced - Test 22 (10 Apr 2024)-img/p2_phy_s1_q01_D.png}{72.3}{0}
\qfigdef{p2_phy_s1_q03_1}{Score Advanced - Test 22 (10 Apr 2024)-img/p2_phy_s1_q03_1.png}{117.1}{0}
\qfigdef{p2_phy_s1_q04_1}{Score Advanced - Test 22 (10 Apr 2024)-img/p2_phy_s1_q04_1.png}{79.1}{0}
\qfigdef{p2_phy_s1_q06_2}{Score Advanced - Test 22 (10 Apr 2024)-img/p2_phy_s1_q06_2.png}{128.1}{0}
\qfigdef{p2_phy_s2_q07_1}{Score Advanced - Test 22 (10 Apr 2024)-img/p2_phy_s2_q07_1.png}{108.3}{0}
\qfigdef{p2_phy_s2_q08_1}{Score Advanced - Test 22 (10 Apr 2024)-img/p2_phy_s2_q08_1.png}{114.8}{0}
\qfigdef{p2_phy_s2_q09_1}{Score Advanced - Test 22 (10 Apr 2024)-img/p2_phy_s2_q09_1.png}{132.2}{0}
\qfigdef{p2_phy_s2_q10_1}{Score Advanced - Test 22 (10 Apr 2024)-img/p2_phy_s2_q10_1.png}{141.5}{0}
\qfigdef{p2_phy_s2_q11_1}{Score Advanced - Test 22 (10 Apr 2024)-img/p2_phy_s2_q11_1.png}{104.9}{0}
\qfigdef{p2_phy_s2_q14_2}{Score Advanced - Test 22 (10 Apr 2024)-img/p2_phy_s2_q14_2.png}{104.6}{0}
\qfigdef{p2_phy_s2_q14_1}{Score Advanced - Test 22 (10 Apr 2024)-img/p2_phy_s2_q14_1.png}{145.7}{0}
\qfigdef{p2_phy_s3_para15_16_1}{Score Advanced - Test 22 (10 Apr 2024)-img/p2_phy_s3_para15_16_1.png}{83.0}{0}
\qfigdef{p2_phy_s3_para17_18_1}{Score Advanced - Test 22 (10 Apr 2024)-img/p2_phy_s3_para17_18_1.png}{92.7}{0}
\qfigdef{p2_phy_s3_para17_18_2}{Score Advanced - Test 22 (10 Apr 2024)-img/p2_phy_s3_para17_18_2.png}{93.4}{0}
\qfigdef{p2_chem_s1_q01_1}{Score Advanced - Test 22 (10 Apr 2024)-img/p2_chem_s1_q01_1.png}{186.6}{0}
\qfigdef{p2_chem_s1_q01_A}{Score Advanced - Test 22 (10 Apr 2024)-img/p2_chem_s1_q01_A.png}{89.4}{0}
\qfigdef{p2_chem_s1_q01_B}{Score Advanced - Test 22 (10 Apr 2024)-img/p2_chem_s1_q01_B.png}{79.3}{0}
\qfigdef{p2_chem_s1_q01_C}{Score Advanced - Test 22 (10 Apr 2024)-img/p2_chem_s1_q01_C.png}{90.1}{0}
\qfigdef{p2_chem_s1_q01_D}{Score Advanced - Test 22 (10 Apr 2024)-img/p2_chem_s1_q01_D.png}{71.7}{0}
\qfigdef{p2_chem_s1_q02_A}{Score Advanced - Test 22 (10 Apr 2024)-img/p2_chem_s1_q02_A.png}{60.8}{0}
\qfigdef{p2_chem_s1_q02_B}{Score Advanced - Test 22 (10 Apr 2024)-img/p2_chem_s1_q02_B.png}{75.2}{0}
\qfigdef{p2_chem_s1_q02_C}{Score Advanced - Test 22 (10 Apr 2024)-img/p2_chem_s1_q02_C.png}{74.2}{0}
\qfigdef{p2_chem_s1_q02_D}{Score Advanced - Test 22 (10 Apr 2024)-img/p2_chem_s1_q02_D.png}{70.8}{0}
\qfigdef{p2_chem_s1_q03_1}{Score Advanced - Test 22 (10 Apr 2024)-img/p2_chem_s1_q03_1.png}{245.2}{0}
\qfigdef{p2_chem_s1_q03_A}{Score Advanced - Test 22 (10 Apr 2024)-img/p2_chem_s1_q03_A.png}{68.0}{0}
\qfigdef{p2_chem_s1_q03_B}{Score Advanced - Test 22 (10 Apr 2024)-img/p2_chem_s1_q03_B.png}{72.9}{0}
\qfigdef{p2_chem_s1_q03_C}{Score Advanced - Test 22 (10 Apr 2024)-img/p2_chem_s1_q03_C.png}{76.1}{0}
\qfigdef{p2_chem_s1_q03_D}{Score Advanced - Test 22 (10 Apr 2024)-img/p2_chem_s1_q03_D.png}{71.0}{0}
\qfigdef{p2_chem_s1_q04_1}{Score Advanced - Test 22 (10 Apr 2024)-img/p2_chem_s1_q04_1.png}{56.5}{0}
\qfigdef{p2_chem_s2_q07_A}{Score Advanced - Test 22 (10 Apr 2024)-img/p2_chem_s2_q07_A.png}{106.8}{0}
\qfigdef{p2_chem_s2_q07_B}{Score Advanced - Test 22 (10 Apr 2024)-img/p2_chem_s2_q07_B.png}{116.1}{0}
\qfigdef{p2_chem_s2_q07_C}{Score Advanced - Test 22 (10 Apr 2024)-img/p2_chem_s2_q07_C.png}{105.6}{0}
\qfigdef{p2_chem_s2_q07_D}{Score Advanced - Test 22 (10 Apr 2024)-img/p2_chem_s2_q07_D.png}{107.3}{0}
\qfigdef{p2_chem_s2_q08_1}{Score Advanced - Test 22 (10 Apr 2024)-img/p2_chem_s2_q08_1.png}{235.8}{0}
\qfigdef{p2_chem_s2_q09_1}{Score Advanced - Test 22 (10 Apr 2024)-img/p2_chem_s2_q09_1.png}{86.3}{0}
\qfigdef{p2_chem_s2_q09_A}{Score Advanced - Test 22 (10 Apr 2024)-img/p2_chem_s2_q09_A.png}{53.2}{0}
\qfigdef{p2_chem_s2_q09_B}{Score Advanced - Test 22 (10 Apr 2024)-img/p2_chem_s2_q09_B.png}{74.2}{0}
\qfigdef{p2_chem_s2_q09_C}{Score Advanced - Test 22 (10 Apr 2024)-img/p2_chem_s2_q09_C.png}{54.1}{0}
\qfigdef{p2_chem_s2_q09_D}{Score Advanced - Test 22 (10 Apr 2024)-img/p2_chem_s2_q09_D.png}{69.3}{0}
\qfigdef{p2_chem_s2_q10_1}{Score Advanced - Test 22 (10 Apr 2024)-img/p2_chem_s2_q10_1.png}{167.3}{0}
\qfigdef{p2_chem_s3_para15_16_1}{Score Advanced - Test 22 (10 Apr 2024)-img/p2_chem_s3_para15_16_1.png}{201.1}{0}
\qfigdef{p2_chem_s3_q15_1}{Score Advanced - Test 22 (10 Apr 2024)-img/p2_chem_s3_q15_1.png}{198.5}{0}
\qfigdef{p2_chem_s3_q15_A}{Score Advanced - Test 22 (10 Apr 2024)-img/p2_chem_s3_q15_A.png}{93.7}{0}
\qfigdef{p2_chem_s3_q15_B}{Score Advanced - Test 22 (10 Apr 2024)-img/p2_chem_s3_q15_B.png}{95.1}{0}
\qfigdef{p2_chem_s3_q15_C}{Score Advanced - Test 22 (10 Apr 2024)-img/p2_chem_s3_q15_C.png}{93.8}{0}
\qfigdef{p2_chem_s3_q15_D}{Score Advanced - Test 22 (10 Apr 2024)-img/p2_chem_s3_q15_D.png}{93.5}{0}
\qfigdef{p2_chem_s3_q16_1}{Score Advanced - Test 22 (10 Apr 2024)-img/p2_chem_s3_q16_1.png}{94.5}{0}
\qfigdef{p2_chem_s3_q16_A}{Score Advanced - Test 22 (10 Apr 2024)-img/p2_chem_s3_q16_A.png}{158.2}{0}
\qfigdef{p2_chem_s3_q16_B}{Score Advanced - Test 22 (10 Apr 2024)-img/p2_chem_s3_q16_B.png}{157.7}{0}
\qfigdef{p2_chem_s3_q16_C}{Score Advanced - Test 22 (10 Apr 2024)-img/p2_chem_s3_q16_C.png}{157.9}{0}
\qfigdef{p2_chem_s3_q16_D}{Score Advanced - Test 22 (10 Apr 2024)-img/p2_chem_s3_q16_D.png}{158.7}{0}
\raggedcolumns
\makeatletter\@beginparpenalty=10000\makeatother
\newcolumntype{L}[1]{>{\raggedright\arraybackslash}p{#1}}
\begin{document}
\begin{center}{\Large\bfseries Score Advanced --- Test 22 (10 Apr 2024)}\end{center}
\begin{multicols}{2}
\raggedcolumns
\papersec{1}
\partsec{1}{PHYSICS}
\sect{SECTION-I (i)}
\qn{1} A plank is held at an angle $\alpha$ to the horizontal on two fixed supports A and B. The plank can slide against the supports (with out friction) because of its weight $Mg$. With what acceleration and in which direction, a man of mass $m$ should move so that the plank does not move?\par
\qfig{p1_phy_s1_q01_1}
\begin{qoptions}
\opt{A}{$g\sin\alpha\left(1+\dfrac mM\right)$ down the incline}
\opt{B}{$g\sin\alpha\left(1+\dfrac Mm\right)$ down the incline}
\opt{C}{$g\sin\alpha\left(1+\dfrac mM\right)$ up the incline}
\opt{D}{$g\sin\alpha\left(1+\dfrac Mm\right)$ up the incline}
\end{qoptions}
\qn{2} With standard resistances 2.0 $\Omega$ and 1.0 $\Omega$ in the gaps of a metre bridge, the balancing length is 66.9 cm from the higher resistance A of the bridge wire. When the resistances are interchanged, the balancing length is 33.5 cm from the same end A. With an unknown resistance $X$ and standard 1.0 $\Omega$ in the gaps, the balancing length is 59.7 cm from the end A nearer to X. Then value of X is
\begin{qoptions}
\opt{A}{0.675 $\Omega$}
\opt{B}{0.670 $\Omega$}
\opt{C}{0.680 $\Omega$}
\opt{D}{1.481 $\Omega$}
\end{qoptions}
\qn{3} A thick hollow sphere with inner radius R and outer radius 3R has a uniform volume density $\rho$ . It has a spherical cavity of radius R as shown. The gravitational field at the centre O' of the cavity is\par
\qfig{p1_phy_s1_q03_1}
\begin{qoptions}
\opt{A}{zero}
\opt{B}{$\dfrac23\pi G\rho R$}
\opt{C}{$\dfrac53\pi G\rho R$}
\opt{D}{$\dfrac73\pi G\rho R$}
\end{qoptions}
\qn{4} Circuit shown in figure is a simple ohm $-$ meter, in which G is a galvanometer, $R_0$ is known resistance and R is the resistance which is to be measured. If A and B are short circuited by a resistancesless wire, galvanometer gives full scale deflection. Then to read the resistance R directly from galvanometer, ohm $-$ meter scale would look like:\par
\qfig{p1_phy_s1_q04_1}
\begin{qoptions}
\opt{A}{\qfigi{p1_phy_s1_q04_A}}
\opt{B}{\qfigi{p1_phy_s1_q04_B}}
\opt{C}{\qfigi{p1_phy_s1_q04_C}}
\opt{D}{\qfigi{p1_phy_s1_q04_D}}
\end{qoptions}
\qn{5} Two identical balls A \& B of mass $m$ each are placed on a fixed wedge as shown in figure. Ball B is kept at rest and it is released just before two balls collides. Ball A rolls down without slipping on inclined plane \& collide elastically with ball B. The kinetic energy of ball A just after the collision with ball B is :\par
\qfig{p1_phy_s1_q05_1}
\begin{qoptions}
\opt{A}{$\dfrac{mgh}{7}$}
\opt{B}{$\dfrac{mgh}{2}$}
\opt{C}{$\dfrac{2mgh}{5}$}
\opt{D}{$\dfrac{7mgh}{5}$}
\end{qoptions}
\sect{SECTION-I (ii)}
\qn{6} An open bottle lies on a precise weighing scale. A fly then enters the bottle and flutters around in futile attempt to escape. Another fly, instead of using its wings to fly, uses a small helium filled balloon fitted to it by a microartist, in order to remain airborne. Choose the correct option(s) from the following statements:
\begin{qoptions}
\opt{A}{reading will more in first case when it flutters}
\opt{B}{reading will more in second case when it is using balloon}
\opt{C}{reading in second cases is higher than reading before the fly entered the bottle.}
\opt{D}{reading will be same in both cases and equal to reading before any fly entered the bottle}
\end{qoptions}
\qn{7} A double conical pendulum consists of two masses, $m$ and $M$, connected by a massless string passing over a frictionless, massless pulley. The entire apparatus rotates clockwise and freely at constant angular speed $\omega$ (rad/s) about the vertical axis (dashed line) passing through centre of pulley as shown. After the system comes in steady state, the length of string on either sides of pulley are $l$ and $L$. Pick up the correct option from the following:\par
\qfig{p1_phy_s2_q07_1}
\begin{qoptions}
\opt{A}{$\dfrac{\cos\phi}{\cos\theta}=\dfrac mM$}
\opt{B}{$\cos\theta=\dfrac{g}{\omega^2L}$}
\opt{C}{$Ml=mL$}
\opt{D}{$\cos\phi=\dfrac{g}{\omega^2l}$}
\end{qoptions}
\qn{8} In case of superposition of waves (at $x=0$)\par
$y_1=2\sin(1026\pi t)$\par
$y_2=2\sin(1014\pi t)$
\begin{qoptions}
\opt{A}{the frequency of resulting wave is 510 Hz}
\opt{B}{the amplitude of resulting wave varies at frequency 3 Hz}
\opt{C}{the frequency of beats is 6 Hz}
\opt{D}{the amplitude of resulting wave varies at frequency 6 Hz}
\end{qoptions}
\qn{9} A,B and C are three large, parallel conducting plates, placed horizontally . A and C are rigidly fixed and earthed. B is given some charge. Under electrostatic and gravitational forces, B may be:\par
\qfig{p1_phy_s2_q09_1}
\begin{qoptions}
\opt{A}{in equilibrium midway between A and C}
\opt{B}{in equilibrium if it is closer to A than to C}
\opt{C}{in equilibrium if it is closer to C than to A}
\opt{D}{B can never be in stable equilibrium}
\end{qoptions}
\qn{10} A particle of specific charge $q/m=-2\times10^{10}$C/kg is travelling with a uniform velocity $v_0\hat{i}$. A uniform steady magnetic field $\vec B$ is suitability applied for a duration of time $\Delta t=\pi\mu$s so that the particle's velocity just after removal of the field becomes $-v_0\hat{i}$. If the motion of the particle is confined tot he XOY plane, $\vec B$ is possibly
\begin{qoptions}
\opt{A}{$50\hat{k}\ \mu$T}
\opt{B}{$-100\hat{k}\ \mu$T}
\opt{C}{$+150\hat{k}\ \mu$T}
\opt{D}{$-250\hat{k}\ \mu$T}
\end{qoptions}
\qn{11} A particle of charge '$q$' and mass '$m$' moving with speed '$u$' along x-axis at origin enters the circular region of a uniform magnetic field $\vec B=B_0\hat{k}$. The particle enters the region from origin at $t=0$ and comes out of the region from point having co-ordinates $(x, y, 0)$ at time $t=t_0$ with speed '$v$'. If the region is gravity free space and the radius of circular region is $R=\dfrac{mu}{qB_0}\sqrt3$ then\par
\qfig{p1_phy_s2_q11_1}
\begin{qoptions}
\opt{A}{$x=\dfrac{mu}{qB_0}\left(\dfrac{\sqrt2-1}{\sqrt2}\right)$, $y=\dfrac{mu}{\sqrt2(qB_0)}$}
\opt{B}{$v=u$}
\opt{C}{$t_0=\dfrac{\pi m}{2qB_0}$}
\opt{D}{$t_0=\dfrac{2\pi m}{3qB_0}$}
\end{qoptions}
\qn{12} A conductor is made of an isotropic material (resistivity $\rho$ ) has rectangular cross$-$section. Horizontal dimension of the rectangle decreases linearly from $2x$ at one end to $x$ at the other end and vertical dimension increases from $y$ to $2y$ as shown in figure. Length of the conductor along the axis is equal to $l$. A battery is connected across the length of this conductor. Then:
\begin{qoptions}
\opt{A}{Resistance of the conductor is equal to $\dfrac{2\rho l}{xy}\log_e(2)$}
\opt{B}{Rate of generation of heat per unit length is maximum at middle cross - section.}
\opt{C}{Drift velocity of conduction electrons is minimum at middle section.}
\opt{D}{At the ends of the conductor, electric field intensity is same.}
\end{qoptions}
\qn{13} A dipole of dipole moment $p$ is placed at a distance $2R$ from the centre of a neutral conducting sphere of radius $R$. The direction of dipole is towards the centre of the sphere. A tangent is drawn from the dipole to the sphere which meets the sphere at point A.
\begin{qoptions}
\opt{A}{The potential at point A is $p/16\pi\varepsilon_0R^2$}
\opt{B}{The potential at point A is $\sqrt3p/24\pi\varepsilon_0R^2$}
\opt{C}{The potential at point A due to induced charges is $\left(3-2\sqrt3\right)p/48\pi\varepsilon_0R^2$}
\opt{D}{The potential at point A due to dipole is $\sqrt3p/24\pi\varepsilon_0R^2$}
\end{qoptions}
\sect{SECTION-III}
\qn{1} Two particles A and B can execute SHM of amplitude `$a$' each and time periods 4s and 6s respectively along the x-axis with origin as their equilibrium position. At time $t=0$, the particle A is at $x=-a$, B is at $x=0$ and both are moving in the same direction. The instant of time at which they collide with each other for the first time is $9/n$ sec. find $n$
\qn{2} A source of sonic oscillations with frequency $f_0=1700$ Hz and a receiver are located at the same point. At the moment $t=0$ the source starts receding from the receiver with constant acceleration $a=10.0$ m/s$^2$ . Assuming the velocity of sound to be equal to $v=340$ m/s . The oscillation frequency registered by the stationary receiver at $t=10.0$ second after the start of motion is $50n$ Hz less than the natural frequency of the source. find $n$.
\qn{3} A 20 gm bullet is fired into a 4 kg wooden block and becomes embedded in it. Knowing that the block and bullet then move up the smooth long inclined (fixed) for 12 mili sec before they come to a stop.\par
The magnitude of the initial velocity of the bullet is $4n$ m/sec, find $n$ (take $g=10$ m/s$^2$)\par
\qfig{p1_phy_s3_q03_1}
\qfig{p1_phy_s3_q03_2}
\qn{4} Two conducting square plates with length $l$ are arranged parallel to each other at a distance $d$ ($\ll l$). The space between the plates is filled with a dielectric of relative permittivity $\varepsilon_r=4$ up to a length $x=l/3$. The upper plate is given a charge $Q$ and the lower plate is given a charge $-Q$ as shown in the figure. If $Q=6\,\mu$C then find how much bound charge (in $\mu$C) is present on the lower surface of dielectric?\par
\qfig{p1_phy_s3_q04_1}
\qn{5} A container filled with air under pressure $P_0$ contains a soap bubble of radius $R$. The air pressure has been reduced to half isothermally and the new radius of the bubble becomes $\dfrac{5R}{4}$. If the surface tension of the soap water solution is $S$, $P_0$ is found to be $\dfrac{16nS}{R}$ SI unit. Find $n$.
\partsec{2}{CHEMISTRY}
\sect{SECTION-I (i)}
\qn{1} The correct statement among the following is
\begin{qoptions}
\opt{A}{When sodium nitrate is heated with sodium metal the products formed are Na$_2$O and N$_2$O}
\opt{B}{When sodium nitrite is heated with sodium metal the products formed are Na$_2$O and N$_2$}
\opt{C}{When sodium nitrate is fused with sodium metal the products formed are NaNO$_2$ and Na$_2$O}
\opt{D}{When sodium nitrite is heated with sodium metal the products formed are Na$_2$O and NO}
\end{qoptions}
\qn{2} A deliquescent white crystalline hydroxide X react with a metal nitrate (Y) to from another hydroxide which decomposes to give insoluble brown precipitate. (X) and (Y) are
\begin{qoptions}
\opt{A}{NaOH, Zn(NO$_3$)$_2$}
\opt{B}{NaOH, AgNO$_3$}
\opt{C}{Ca(OH)$_2$, Hg(NO$_3$)$_2$}
\opt{D}{NaOH, Al(NO$_3$)$_3$}
\end{qoptions}
\qn{3} \qfig{p1_chem_s1_q03_1}
\begin{qoptions}
\opt{A}{\qfigi{p1_chem_s1_q03_A}}
\opt{B}{\qfigi{p1_chem_s1_q03_B}}
\opt{C}{\qfigi{p1_chem_s1_q03_C}}
\opt{D}{\qfigi{p1_chem_s1_q03_D}}
\end{qoptions}
\qn{4} \qfig{p1_chem_s1_q04_1}
C is:
\begin{qoptions}
\opt{A}{\qfigi{p1_chem_s1_q04_A}}
\opt{B}{\qfigi{p1_chem_s1_q04_B}}
\opt{C}{\qfigi{p1_chem_s1_q04_C}}
\opt{D}{\qfigi{p1_chem_s1_q04_D}}
\end{qoptions}
\qn{5} Which of the following does not exist as a Zwitter ion?
\begin{qoptions}
\opt{A}{Glycine}
\opt{B}{Para Sulphanilic acid}
\opt{C}{\qfigi{p1_chem_s1_q05_C}}
\opt{D}{Alanine}
\end{qoptions}
\sect{SECTION-I (ii)}
\qn{6} When KF is dissolved in HF, a new species KHF$_2$ is formed. The type of bond present in KHF$_2$ is/are
\begin{qoptions}
\opt{A}{Ionic bond}
\opt{B}{Covalent bond}
\opt{C}{Coordinate covalent bond}
\opt{D}{Hydrogen bond}
\end{qoptions}
\qn{7} HNO$_3$ + P$_4$O$_{10}$ $\rightarrow$ A + HPO$_3$\par
A + Na $\rightarrow$ B + C\par
Cu + HNO$_3$ (Conc) $\rightarrow$ Cu(NO$_3$)$_2$ + C + H$_2$O\par
C + NaOH $\rightarrow$ B + D\par
Select the correct statements
\begin{qoptions}
\opt{A}{B acts as a base in molten A}
\opt{B}{D acts as an acid in molten A}
\opt{C}{C has very less tendency toward dimerisation at room temperature}
\opt{D}{All the atoms in A are in same plane}
\end{qoptions}
\qn{8} Which of the following ions do not contain electron in ns (4s) orbital
\begin{qoptions}
\opt{A}{Mn$^{+}$}
\opt{B}{Cr$^{+}$}
\opt{C}{V$^{+}$}
\opt{D}{Fe$^{+}$}
\end{qoptions}
\qn{9} Which of the following statements are correct
\begin{qoptions}
\opt{A}{The transition metal ion is more polarizing compared to S-block metal ions of comparable size for a particular oxidation state}
\opt{B}{The solubility of silver halides is in the order AgF $>$ AgCl $>$ AgBr $>$ AgI and the decrease in solubility is due to increased covalent character by the increased polarisability of halide ions}
\opt{C}{LiCl is soluble in organic solvents while NaCl is insoluble in such solvents due to less covalent character}
\opt{D}{Thermal stability of oxosalts decreases with decrease in the size of cation}
\end{qoptions}
\qn{10} The energy profiles for the two possible mechanisms of the following reaction\par
\qfig{p1_chem_s2_q10_1}
Are given below\par
\qfig{p1_chem_s2_q10_2}
Consider the following substrates (I to IV) for the above reaction\par
(I) R$_1$ = R$_2$ = R$_3$ = H\par
(II) R$_1$ = C$_2$H$_5$ ; R$_2$ = R$_3$ = H\par
(III) R$_1$ = C$_2$H$_5$ ; R$_2$ = CH$_3$ ; R$_3$ = H\par
(IV) R$_1$ = C$_2$H$_5$ ; R$_2$ = R$_3$ = CH$_3$\par
Choose the correct statements
\begin{qoptions}
\opt{A}{The reaction mechanism changes from the one shown in figure Y to that in figure X, as one changes the substrate from I to IV}
\opt{B}{Considering the mechanism corresponding to figure X for substrates I to IV, the correct order of $\Delta$E is $\Delta E_{I}<\Delta E_{II}<\Delta E_{III}<\Delta E_{IV}$}
\opt{C}{The reaction for substrate I follows 2$^{nd}$ order kinetics and that for IV follows 1$^{st}$ order kinetics}
\opt{D}{If CN$^{-}$ is taken in large excess, it is difficult to ascertain if the reaction follows mechanism corresponding to energy profile in figure X or that in figure Y}
\end{qoptions}
\qn{11} The nitration of toluene takes place at the ortho, para and meta position, giving ortho-, para-, and meta nitrotoluenes, respectively. Assume that the meta product is formed in negligible amount and that the ortho and para substitutions are irreversible first $-$ order reactions. Let the initial toluene concentration be [A]$_0$ and k$_1$ and k$_2$ be the ortho- and para $-$ rate constants, respectively. Choose the correct statements
\begin{qoptions}
\opt{A}{The rate of formation of the ortho substituted product is given by $k_1[A]\,e^{-(k_1+k_2)t}$, where t is the time.}
\opt{B}{If the rate constant for the ortho product formation is twice that for the para product formation, then [O-nitrotoluene]=2 [p-nitrotoluene], at any time}
\opt{C}{Nitration at the ortho $-$ and para $-$ positions is favorable compared to that at the meta position as the ortho and para $-$ intermediates are more stabilized than the meta $-$ intermediate}
\opt{D}{If the methyl group of toluene is replaced by tertiary butyl group, more ortho product is expected}
\end{qoptions}
\qn{12} The following four compounds are subjected to different reactions\par
\qfig{p1_chem_s2_q12_1}
Choose the correct statements for these reactions
\begin{qoptions}
\opt{A}{Compound I on reaction with bromine and NaOH gives the same compound as that obtained by complete reduction of III.}
\opt{B}{The compound that can be reduced by hydrogen under pressure in the presence of Raney Ni is IV.}
\opt{C}{Compound III on reaction with LiAlH$_4$ gives a compound which reacts with bromine under suitable condition.}
\opt{D}{The compound obtained by the reaction of LiAlH$_4$ with II is same as that obtained on heating IV with dilute H$_2$SO$_4$ at high temperature.}
\end{qoptions}
\qn{13} Which of the following will produce aromatic compound as product?
\begin{qoptions}
\opt{A}{CH$_3-$C$\equiv$CH $\xrightarrow[\Delta]{\text{Red hot Fe tube}}$}
\opt{B}{$n$C$_6$H$_{14}$ $\xrightarrow[\text{Cr}_2\text{O}_3/600^\circ\text{C}]{\text{Al}_2\text{O}_3}$ ......}
\opt{C}{\qfigi{p1_chem_s2_q13_1}}
\opt{D}{\qfigi{p1_chem_s2_q13_D}}
\end{qoptions}
\sect{SECTION-III}
\qn{1} Chlorine oxidises hypo to a product (X). The difference in the sum of the oxidation numbers of sulphur in hypo and sulphur in product (X) is
\qn{2} A sulphate of a metal (M) on heating evolves two gases (B) and (C) and an oxide (D). Gas (B), turns acidified dichromate paper green while gas (C) forms a trimer.\par
Compound (D) forms a Lewis acid (E) with HCl (g). If the number of S $-$ O $-$ S bonds in trimer of gas (C) be `n' and the number of equivalent M $-$ Cl bonds in one molecule of (E) in vapour phase be `m' then (m+n) is
\qn{3} Total number of possible isomers of the complex [Cr(en)$_3$] [Cr(OX)$_3$] is
\qn{4} How many of the following is/are condensation polymer? Teflon, Nylon 66, PVC, Buna-S, Bakelite, Terylene, Neoprene, Silicones, HDPE.
\qn{5} How many are the correct statements ?\par
1) CH (CN)$_3$ is more acidic than Picric acid\par
2) Tri ethyl amine is more basic than \qfigi{p1_chem_s3_q05_1} (Quinuclidine)\par
3) In propene the eclipsed conformation lies lesser potential energy than bisected Conformation\par
4) NH$_4$Cl, naphthalene, I$_2$, indigo are all compounds which sublime easily
\partsec{3}{MATHEMATICS}
\sect{SECTION-I (i)}
\qn{1} An ellipse in the first quadrant is tangent to co-ordinate axes. If one of focus is F$_1$ (3,7), and the other focus is F$_2$ (d,7), then the value of $\sqrt{3d}$ is equal to
\begin{qoptions}
\opt{A}{5}
\opt{B}{7}
\opt{C}{14}
\opt{D}{49}
\end{qoptions}
\qn{2} Let f : $\mathbb R\to\mathbb R^{+}$ be a continuous function such that $f(x)=2(1+x^2)\left(1+\displaystyle\int_0^x\dfrac{f(t)}{1+t^2}\,dt\right),\ \forall x\in\mathbb R$. Then
\begin{qoptions}
\opt{A}{f(2)=6e$^4$}
\opt{B}{f(1)=4e}
\opt{C}{f(2)=10e$^4$}
\opt{D}{f(3)=8e$^3$}
\end{qoptions}
\qn{3} $\displaystyle\lim_{n\to\infty}\dfrac1{\sqrt n}\sum_{r=1}^{n}\dfrac1{\sqrt{r+\sqrt{r^2-1}}}$ is
\begin{qoptions}
\opt{A}{0}
\opt{B}{$\dfrac12$}
\opt{C}{1}
\opt{D}{$\dfrac1{\sqrt2}$}
\end{qoptions}
\qn{4} Let g(x) = max$|y^2-xy|$, y $\in$ [0, 1], then the absolute minimum value of g(x) is
\begin{qoptions}
\opt{A}{1}
\opt{B}{0}
\opt{C}{$3-2\sqrt2$}
\opt{D}{$2\sqrt2-2$}
\end{qoptions}
\qn{5} $\displaystyle\lim_{n\to\infty}\left(\binom n0\binom n1\binom n2\cdots\binom nn\right)^{\frac1{n(n+1)}}$ is
\begin{qoptions}
\opt{A}{$\dfrac1{\sqrt2}$}
\opt{B}{$\sqrt2$}
\opt{C}{$\dfrac1{\sqrt e}$}
\opt{D}{$\sqrt e$}
\end{qoptions}
\sect{SECTION-I (ii)}
\qn{6} Let Z be a variable complex number satisfying $\dfrac{2z-\bar z}{3}=a-\dfrac13+ib$ where a, b are real parameters such that $a^2-b^2=1$. If locus of z is a curve `C', then
\begin{qoptions}
\opt{A}{director circle of the curve `C' is $(x+1)^2+y^2=10$}
\opt{B}{area of the triangle formed by the straight line $x-3y+1=0$ and tangent to the curve `C' at one of its vertex and the x-axis is $\dfrac32$ sq. units}
\opt{C}{Distance between directrices of the curve `C' is $\dfrac{18}{\sqrt{10}}$}
\opt{D}{Equation of the circle passing through the points (2,1) and both foci of the curve `C' is $x^2+y^2+2x-9=0$}
\end{qoptions}
\qn{7} If M$_1$ and M$_2$ are the feet of perpendiculars from foci F$_1$ and F$_2$ of the ellipse $\dfrac{x^2}{64}+\dfrac{y^2}{25}=1$ on the tangent at any point P of the ellipse respectively. If M$_1$F$_1$ = p and M$_2$F$_2$ = q then
\begin{qoptions}
\opt{A}{min $(p^2+q^2)=50$}
\opt{B}{min $(p^2+q^2)=25$}
\opt{C}{min $(p+q)=10$}
\opt{D}{min $(p+q)=50$}
\end{qoptions}
\qn{8} Let AB : y = x be a chord to the parabola $y^2+4x=4$ and P be any variable point on the ellipse $x^2+9(y-4)^2-9=0$. If $\Delta_1$ and $\Delta_2$ be the minimum and maximum area of triangle ABP respectively then
\begin{qoptions}
\opt{A}{$\Delta_2=2\sqrt2\left(4+\sqrt{10}\right)$}
\opt{B}{$\Delta_2-\Delta_1=8\sqrt5$}
\opt{C}{$\dfrac{\Delta_2}{\Delta_1}=\dfrac{13+4\sqrt{10}}{3}$}
\opt{D}{$\Delta_2-\Delta_1=5\sqrt5$}
\end{qoptions}
\qn{9} $f(x)=\displaystyle\lim_{n\to\infty}\sum_{r=1}^{n}3^{r-1}\sin^3\dfrac{x}{3^r}$, then
\begin{qoptions}
\opt{A}{$f\left(\dfrac\pi2\right)=\dfrac\pi8-\dfrac14$}
\opt{B}{$f\left(\dfrac\pi2\right)=\dfrac\pi8$}
\opt{C}{$f\left(\dfrac{3\pi}2\right)=\dfrac{3\pi}8$}
\opt{D}{$f\left(\dfrac{3\pi}2\right)=\dfrac{3\pi}8+\dfrac14$}
\end{qoptions}
\qn{10} Let $f(x,y,z)=xyz$ and $A=\{(x,y,z)\in\mathbb R^3:x^2+y^2+z^2=1,\ x+y+z=1\}$. Then over set A
\begin{qoptions}
\opt{A}{global maximum of f (x, y, z) is 0}
\opt{B}{global maximum of f (x, y, z) is $-\dfrac4{27}$}
\opt{C}{global maximum of x is 1}
\opt{D}{global minimum of z is $-\dfrac13$}
\end{qoptions}
\qn{11} Let m, n $\in\mathbb N$ such that n $>$ m $>$ 2. Consider a cubic equation $mx^3+nx^2-1=0$ with $\alpha,\beta,\gamma$ as its roots. Which of the following statements are true ?
\begin{qoptions}
\opt{A}{All three roots $\alpha,\beta,\gamma$ are real.}
\opt{B}{Two roots are negative and one is positive}
\opt{C}{If $\alpha<\beta<\gamma$ then $|\gamma|<|\beta|<|\alpha|$}
\opt{D}{If $\alpha<\beta<\gamma$ then $|\beta|<|\gamma|<|\alpha|$}
\end{qoptions}
\qn{12} Let f be a differentiable function on [0, 1] such that f(0) = 0 and f(1) = 1. Then
\begin{qoptions}
\opt{A}{There exist $x_1,x_2\in(0,1)$ and $x_1\ne x_2$ such that $f'(x_1)+f'(x_2)=2$}
\opt{B}{There exist $x_1,x_2\in(0,1)$ and $x_1\ne x_2$ such that $\dfrac1{f'(x_1)}+\dfrac1{f'(x_2)}=2$}
\opt{C}{Corresponding to arbitrary positive real numbers $k_1$ and $k_2$, there exist $x_1,x_2\in(0,1)$ and $x_1\ne x_2$ such that $\dfrac{k_1}{f'(x_1)}+\dfrac{k_2}{f'(x_2)}=k_1+k_2$.}
\opt{D}{There exists n distinct real numbers $x_1,x_2,\dots,x_n\in(0,1)$ such that $f'(x_1)+f'(x_2)+\dots+f'(x_n)=n$}
\end{qoptions}
\qn{13} Consider two curves namely $C_1: y=a^x$ and $C_2: x=a^y$ ; $a>0$ then
\begin{qoptions}
\opt{A}{Number of times $C_1$ meets with $C_2$ is 3 for $a\in(0,e^{-e})$}
\opt{B}{Number of times $C_1$ meets with $C_2$ in 1 for $a\in\left[e^{-e},1\right]\cup\left\{e^{e^{-1}}\right\}$}
\opt{C}{Number of times $C_1$ meets with $C_2$ is 2 for $a\in\left(1,e^{e^{-1}}\right)$}
\opt{D}{Number of times $C_1$ meets with $C_2$ is 0 for $a\in\left(e^{e^{-1}},\infty\right)$}
\end{qoptions}
\sect{SECTION-III}
\qn{1} For $a,b\in\mathbb R^{+}$, $\displaystyle\int_a^b\dfrac{e^{x/a}-e^{b/x}}{x}\,dx$ is
\qn{2} Let $A=\begin{bmatrix}a&b\\c&d\end{bmatrix}$, $B=\begin{bmatrix}p\\q\end{bmatrix}\neq\begin{bmatrix}0\\0\end{bmatrix}$ such that AB = B and tr.(A) = 50 and ad $-$ bc is k then $\lfloor\sqrt k\rfloor$ is
\qn{3} Let $A=\begin{bmatrix}2&1&0\\0&2&1\\0&0&2\end{bmatrix}$. If sum of all elements of $A^{10}$ is `k' then $\left\lfloor\dfrac{k}{2560}\right\rfloor$ is
\qn{4} Number of terms which are odd in the sequence $a_0,a_1,a_2,\dots,a_{2017}$ where $a_n=\dbinom{2017}{n}$ is $\dfrac k2$. Then $\lfloor\sqrt k\rfloor$ is
\qn{5} Number of non constant differentiable function $f:\mathbb R\to\mathbb R^{+}\cup\{0\}$, such that $f(f(f(x)))=f(x),\forall x\in\mathbb R$ is
\raggedcolumns
\papersec{2}
\partsec{1}{PHYSICS}
\sect{SECTION-I (i)}
\qn{1} A small particle of mass m is released from rest from point A inside a fixed smooth hemispherical bowl as shown in the figure. The ratio (r) of magnitude of centripetal force and normal reaction on the particle at any point varies with $\theta$ as\par
\qfig{p2_phy_s1_q01_1}
\begin{qoptions}
\opt{A}{\qfigi{p2_phy_s1_q01_A}}
\opt{B}{\qfigi{p2_phy_s1_q01_B}}
\opt{C}{\qfigi{p2_phy_s1_q01_C}}
\opt{D}{\qfigi{p2_phy_s1_q01_D}}
\end{qoptions}
\qn{2} A simple pendulum consists of a bob of mass m and thread of length $\ell$. In an experimental demonstration, your physics teacher pulls the bob aside until the thread makes an angle $\theta_0=\left(\dfrac{2}{\sqrt3}-1\right)$ radians with the vertical. Now he pushes the bob towards the equilibrium position. While pushing, he continuously applies a force F = mg along the circular arc until the thread comes in vertical position and thereafter he releases the bob. For what approximate length of time $t_0$ did he push the bob ?
\begin{qoptions}
\opt{A}{$\dfrac\pi6\sqrt{\dfrac\ell g}$}
\opt{B}{$\pi\sqrt{\dfrac\ell g}$}
\opt{C}{$\dfrac\pi3\sqrt{\dfrac\ell g}$}
\opt{D}{$\dfrac{2\pi}3\sqrt{\dfrac\ell g}$}
\end{qoptions}
\qn{3} An engineer designs a robot that can climb stairs. Time T in which battery used in the robot is fully discharged depends on its constant speed v relative to the stairs. This dependence is shown in the graph which follows y$-$0.04 = $x^2$ . Determine the maximum length of the staircase that the robot can climb with constant velocity.\par
\qfig{p2_phy_s1_q03_1}
\begin{qoptions}
\opt{A}{2.5m}
\opt{B}{5m}
\opt{C}{7.5m}
\opt{D}{1.25m}
\end{qoptions}
\qn{4} From a box of balls stacked as shown, balls are removed at equal intervals of time with a speed v as shown. Average velocity of centre of mass of remaining balls in the box is\par
\qfig{p2_phy_s1_q04_1}
\begin{qoptions}
\opt{A}{zero}
\opt{B}{upward}
\opt{C}{downward}
\opt{D}{cannot be determined}
\end{qoptions}
\qn{5} There are 2016 points in space. Each pair of points is connected by a resistance of 1$\Omega$. The equivalent resistance between any two points is
\begin{qoptions}
\opt{A}{$\dfrac1{2016}\Omega$}
\opt{B}{$\dfrac2{2016}\Omega$}
\opt{C}{$\dfrac1{2015}\Omega$}
\opt{D}{$\dfrac2{2015}\Omega$}
\end{qoptions}
\qn{6} A part of a conductor carrying current I is in the form of a semicircle AB of radius R and lying in xy-plane with its centre at origin as shown in figure. Which of the following is the magnitude of $\oint\vec B\cdot d\vec l$ for the circular loop along $3x^2+3z^2=R^2$ in xz-plane due to current in curve AB?\par
\qfig{p2_phy_s1_q06_2}
\begin{qoptions}
\opt{A}{$\dfrac{\sqrt3}{2}\mu_0 I$}
\opt{B}{$\left(\sqrt3+1\right)\dfrac{\mu_0 I}{2}$}
\opt{C}{$\left(2-\sqrt3\right)\dfrac{\mu_0 I}{2}$}
\opt{D}{Zero}
\end{qoptions}
\sect{SECTION-I (ii)}
\qn{7} The figure shows a potentiometer arrangement. D is the driving cell and C is the cell whose emf is to be determined. AB is the potentiometer wire and G is galvanometer. J is sliding contact which can touch any point on AB. Which of the following are essential conditions for obtaining balance ?\par
\qfig{p2_phy_s2_q07_1}
\begin{qoptions}
\opt{A}{The emf of D must be greater than the emf of C}
\opt{B}{Either the positive terminals of both D and C or the negative terminals of both D and C must be joined to A}
\opt{C}{The positive terminals of D and C must be joined}
\opt{D}{The resistance of G must be less than resistance of AB}
\end{qoptions}
\qn{8} A small disc of mass m is attached to the end of a light inextensible cord, which passes through a hole in a frictionless horizontal tabletop. At the other end of the cord, a weight of mass M is attached. Initially disc is moving on a circle of radius R with angular velocity $\omega$. If the hanging weight is pulled slightly down by $x_0$ ($\ll$R) and then released,\par
\qfig{p2_phy_s2_q08_1}
\begin{qoptions}
\opt{A}{The net force on m at an instant when M moves up by $\dfrac{x_0}2$ after the release is $Mg+\dfrac{3x_0mM\omega^2}{2(M+m)}$}
\opt{B}{The net force on M at an instant when M moves up by $\dfrac{x_0}2$ after the release is $Mg+\dfrac{3x_0mM\omega^2}{2(M+m)}$}
\opt{C}{The net force on M at an instant when M moves up by $\dfrac{x_0}2$ after the release is $\dfrac{3x_0mM\omega^2}{2(M+m)}$}
\opt{D}{Acceleration of m at an instant when M moves up by $\dfrac{x_0}2$ after the release is $\dfrac{3x_0m\omega^2}{2(M+m)}$}
\end{qoptions}
\qn{9} Avantika fought her kidnappers with all her might, but to no avail. In her tussle, two flowers (at two different instants) from her bracelet fell in to the river as the boat fled with a constant velocity towards an unknown destination. Probably, Admiral Kattappa was too late, or was he? All he could see in the river were two flowers that once decorated the delicate wrist. The co-ordinates of the flowers were (3m, 3m) and (7m, 6m). Given the river flows from right to left at 3m/s and the speed of the boat with respect to water was 5m/s, in which possible direction(s) w.r.t a fixed xy-axis shown in figure did the kidnappers flee? (Neglect dimensions of the boat)\par
\qfig{p2_phy_s2_q09_1}
\begin{qoptions}
\opt{A}{$\hat i+3\hat j$}
\opt{B}{$-3\hat i+\hat j$}
\opt{C}{$-7\hat i-3\hat j$}
\opt{D}{$\hat i-3\hat j$}
\end{qoptions}
\qn{10} A particle of mass m = 2kg with potential energy shown in the graph is moving towards positive x-axis with a speed 1m/s at x = 1m.\par
\qfig{p2_phy_s2_q10_1}
\begin{qoptions}
\opt{A}{Particle reverses its direction of motion at x = 2m}
\opt{B}{Particle reverses its direction of motion at x = 6m}
\opt{C}{Force acting on the particle at the instant of its direction of reversal is 1N towards negative x direction}
\opt{D}{Maximum speed of the particle before it reverses its direction is 3m/s}
\end{qoptions}
\qn{11} The cable strung between two neighbouring electricity posts sags a little. The mass per unit length of the cable is $\lambda$, the distance between the posts is L, and the `maximum sag' of the cable is d(d $\ll$ L). Due to relatively small sag the tension in the cable can be assumed to be uniform. The cable obeys Hooke's law. Choose the correct statement(s) from the following:\par
\qfig{p2_phy_s2_q11_1}
\begin{qoptions}
\opt{A}{The approximate value of tension is $\dfrac{\lambda L^2g}{4d}$}
\opt{B}{The approximate value of tension is $\dfrac{\lambda L^2g}{8d}$}
\opt{C}{The appropriate formula for calculation of elastic strain energy stored in the cable is $\dfrac12$ x stress x strain x volume}
\opt{D}{The appropriate formula for calculation of elastic strain energy stored in the cable is stress x strain x volume}
\end{qoptions}
\qn{12} The diameter of a cylinder is measured using a vernier calipers with no zero error. While measuring diameter of a cylinder, it is found that the zero of vernier scale lies between 5.10 cm and 5.15 cm of the main scale. Vernier scale has 50 divisions whose total length is equivalent to 2.45 cm. The closest division of the vernier scale that almost coincides with a main scale division is 24$^{th}$ VSD but this 24$^{th}$ VSD actually falls slightly to the right of the main scale division. The diameter of cylinder may be (each MSD = 0.05cm)
\begin{qoptions}
\opt{A}{51.238 mm}
\opt{B}{51.236 mm}
\opt{C}{51.242 mm}
\opt{D}{51.246 mm}
\end{qoptions}
\qn{13} A submarine travels horizontally with a uniform speed `u' inside a calm sea. Then the flow around the submarine is
\begin{qoptions}
\opt{A}{unsteady as seen by a stationary fish in the sea}
\opt{B}{steady as seen by a crab moving with same velocity as that of the marine}
\opt{C}{unsteady as seen by the captain inside the marine}
\opt{D}{steady as seen by a person moving down in the sea}
\end{qoptions}
\qn{14} A source `S' emitting sound of frequency f moves uniformly along a straight line with a velocity u as shown. The velocity of sound in the medium is v($>$u). Two detectors D$_1$ and D$_2$ are placed on the line perpendicular to line of motion of S. Point P represents the point of closest approach of S to the detectors. The frequency detected ($f_D$) by D$_1$ and D$_2$ versus time t graph is shown on the same graph paper.\par
\qfig{p2_phy_s2_q14_2}
\qfig{p2_phy_s2_q14_1}
Choose the incorrect option(s)
\begin{qoptions}
\opt{A}{Graph-1 belongs to detector D$_1$}
\opt{B}{t = $t_0$ is the instant when source S is at P}
\opt{C}{The value of $f_0$ is f}
\opt{D}{The value of $f_0$ is $\left(\dfrac{v^2}{v^2-u^2}\right)$f}
\end{qoptions}
\sect{SECTION-I (iii)}
\paraq{15}{16}
On an air-hockey table, there are N identical small discs lying equally spaced around a semicircle; the total mass of the discs is M. Another small disc, D, of mass m, travelling in a direction perpendicular to the closing diameter of the semicircle, strikes the first of the stationary discs. By some miracle, it subsequently bounces off all of the other N $-$ 1 discs in turn, after which it is travelling in a direction directly opposed to that of its initial motion. All the collisions are perfectly elastic, and friction is everywhere negligible.\par
\qfig{p2_phy_s3_para15_16_1}
\qn{15} In the limiting case of N $\to\infty$, what is the minimal value of the mass ratio M/m for such a miracle to be possible?
\begin{qoptions}
\opt{A}{$\pi$}
\opt{B}{$e^{\pi/2}$}
\opt{C}{4 ln 2}
\opt{D}{2}
\end{qoptions}
\qn{16} When the mass ratio has the critical value found in above question what is the ratio of the final and initial speeds of D?
\begin{qoptions}
\opt{A}{1 ln 2}
\opt{B}{$e^{-\pi}$}
\opt{C}{1}
\opt{D}{$\dfrac12$}
\end{qoptions}
\paraq{17}{18}
Consider first a metal sphere of radius R and total charge Q, which (because of symmetry) is distributed uniformly on its surface, i.e. the surface charge density is $\sigma_0=\dfrac{Q}{4\pi R^2}$ = constant.\par
The total field effects of the charges on the small surface pieces opposite each other cancel out in pairs, so the net electric field of all the charges on the sphere is zero at any point P inside the sphere. The same is shown in figure-1.\par
\qfig{p2_phy_s3_para17_18_1}
\qfig{p2_phy_s3_para17_18_2}
Now let us move on to the case of the electrically charged metal disc. We show that its charge distribution can be found from the uniform surface density of the charged metal sphere. Let us (notionally) `glue into position' the charges on the surface of the sphere, and then project the points of the spherical surface perpendicularly onto a plane that contains the centre of the sphere and point P investigated earlier.\par
As the charges cannot move, the charge $\Delta Q_1$, originally on the surface area $\Delta A_1$, when projected onto the plane, occupies a somewhat smaller area $\Delta A_1^{*}$, and so the surface charge density (with an unchanged amount of charge) increases. In addition, its distance from the point P also changes, from the original value of $r_1$ to the smaller one of $d_1$. The corresponding thing happens with the charge $\Delta Q_2$ originally located on the opposite side of the sphere.
\qn{17} Consider the following statements about figure-1 and figure-2\par
(i) The electric field produced at P by the two charged pieces (projected onto the plane and still glued to the disc) in the figure-2 is zero.\par
(ii) $\dfrac{\Delta A_1}{\Delta A_2}=\dfrac{r_1}{r_2}$\par
(iii) $\dfrac{\Delta A_1}{\Delta A_2}=\dfrac{d_1^2}{d_2^2}$\par
(iv) If a metal disc is charged by Q, the charge distribution on it is the same way as it would be if that on a uniformly charged spherical surface were projected perpendicular on it as shown in figure-2.\par
The correct statement(s) is/are
\begin{qoptions}
\opt{A}{(i), (ii), (iii)}
\opt{B}{(i), (iii), (iv)}
\opt{C}{(i), (ii), (iv)}
\opt{D}{(i), (ii), (iii), (iv)}
\end{qoptions}
\qn{18} A large thin metal disc of radius R is charged and carries a total electric charge of Q. The charge density $\sigma$ on the disc at distance r from centre far away from the rim of the disc is (Neglect edge effects)
\begin{qoptions}
\opt{A}{$\dfrac{Q}{2\pi R^2}\dfrac{R}{\sqrt{R^2-r^2}}$}
\opt{B}{$\dfrac{Q}{4\pi r^2}\dfrac{R}{\sqrt{R^2+r^2}}$}
\opt{C}{$\dfrac{Q}{4\pi R^2}\dfrac{R}{\sqrt{R^2-r^2}}$}
\opt{D}{$\dfrac{Q}{2\pi R^2}\dfrac{R}{\sqrt{R^2+r^2}}$}
\end{qoptions}
\raggedcolumns
\partsec{2}{CHEMISTRY}
\sect{SECTION-I (i)}
\qn{1} Identify the major product in the following reaction?\par
\qfig{p2_chem_s1_q01_1}
\begin{qoptions}
\opt{A}{\qfigi{p2_chem_s1_q01_A}}
\opt{B}{\qfigi{p2_chem_s1_q01_B}}
\opt{C}{\qfigi{p2_chem_s1_q01_C}}
\opt{D}{\qfigi{p2_chem_s1_q01_D}}
\end{qoptions}
\qn{2} What is the structure of the $\gamma$ -lactone formed by iodolactonisation of 4-pentynoic acid?\par
CH $\equiv$ C $-$ CH$_2$ $-$ CH$_2$ $-$ COOH $\xrightarrow[\text{KHCO}_3,\text{H}_2\text{O}]{\text{I}_2.\text{KI}}$ Major product
\begin{qoptions}
\opt{A}{\qfigi{p2_chem_s1_q02_A}}
\opt{B}{\qfigi{p2_chem_s1_q02_B}}
\opt{C}{\qfigi{p2_chem_s1_q02_C}}
\opt{D}{\qfigi{p2_chem_s1_q02_D}}
\end{qoptions}
\qn{3} A key step in the synthesis of an important class of lipids prostaglandins involves the following sequence of reactions. Identify the reactant `x'?\par
\qfig{p2_chem_s1_q03_1}
\begin{qoptions}
\opt{A}{\qfigi{p2_chem_s1_q03_A}}
\opt{B}{\qfigi{p2_chem_s1_q03_B}}
\opt{C}{\qfigi{p2_chem_s1_q03_C}}
\opt{D}{\qfigi{p2_chem_s1_q03_D}}
\end{qoptions}
\qn{4} (B) $\xleftarrow{\text{0.01M NaOH}}$ \qfigi{p2_chem_s1_q04_1} $\xrightarrow{\text{4M NaOH}}$ (A)\par
Which of the following is correct regarding compounds (A) and (B)
\begin{qoptions}
\opt{A}{(A) and (B) are super imposable mirror images}
\opt{B}{The configuration of (A) is `R' and (B) is `S'}
\opt{C}{(A) and (B) are diastereomers}
\opt{D}{(A) is formed with inversion of configuration \& (B) with retention of configuration}
\end{qoptions}
\qn{5} Basing on stability of the lower species to which they are reduced, identify the correct order of oxidising power.
\begin{qoptions}
\opt{A}{$\text{MnO}_4^{-}<\text{VO}_2^{+}<\text{Cr}_2\text{O}_7^{2-}$}
\opt{B}{$\text{Cr}_2\text{O}_7^{2-}<\text{VO}_2^{+}<\text{MnO}_4^{-}$}
\opt{C}{$\text{VO}_2^{+}<\text{MnO}_4^{-}<\text{Cr}_2\text{O}_7^{2-}$}
\opt{D}{$\text{VO}_2^{+}<\text{Cr}_2\text{O}_7^{2-}<\text{MnO}_4^{-}$}
\end{qoptions}
\qn{6} Quartz + Carbon $\xrightarrow[\Delta]{\text{Cl}_2}$ A$_{(\ell)}$ + B$_{(g)}$\par
A + CH$_3$MgBr $\rightarrow$ $\xrightarrow{\text{H}_3\text{O}^{+}}$ C\par
C $\xrightarrow{\Delta}$ D. Identify the correct statement(s) regarding the reaction sequence.
\begin{qoptions}
\opt{A}{B is formed when acetic acid is heated with conc. H$_2$SO$_4$}
\opt{B}{D is a chain silicone}
\opt{C}{D can also be prepared by heating orthosilicic acid.}
\opt{D}{D is a cross linked silicone polymer}
\end{qoptions}
\sect{SECTION-I (ii)}
\qn{7} Identify the pair of acids in which later is more acidic than former?
\begin{qoptions}
\opt{A}{\qfigi{p2_chem_s2_q07_A}}
\opt{B}{\qfigi{p2_chem_s2_q07_B}}
\opt{C}{\qfigi{p2_chem_s2_q07_C}}
\opt{D}{\qfigi{p2_chem_s2_q07_D}}
\end{qoptions}
\qn{8} When the compounds of the following type are allowed to stand in pentane, they are converted into a constitutional isomer. Which of the following is/are correct about the constitutional isomers (A) and (B)?\par
\qfig{p2_chem_s2_q08_1}
\begin{qoptions}
\opt{A}{(A) \& (B) are functional isomers}
\opt{B}{Both give same products on hydrolysis}
\opt{C}{Hydrolysis products can be separated by dil.HCl}
\opt{D}{Hydrolysis products can be separated by dil.NaOH}
\end{qoptions}
\qn{9} What is/are the possible product/s in the following dehydration reaction?\par
\qfig{p2_chem_s2_q09_1}
\begin{qoptions}
\opt{A}{\qfigi{p2_chem_s2_q09_A}}
\opt{B}{\qfigi{p2_chem_s2_q09_B}}
\opt{C}{\qfigi{p2_chem_s2_q09_C}}
\opt{D}{\qfigi{p2_chem_s2_q09_D}}
\end{qoptions}
\qn{10} Which of the following statement/s is/are correct about compounds (A) and (B)?\par
\qfig{p2_chem_s2_q10_1}
\begin{qoptions}
\opt{A}{The rate of hydrogenation of (A) is more than (B) under similar conditions}
\opt{B}{The compounds (A) and (B) are stereo isomers.}
\opt{C}{Catalytic hydrogenation of (A) or (B) gives two diastereomers in unequal ratio.}
\opt{D}{The pair of diastereomers can be resolved in to optically active compounds.}
\end{qoptions}
\qn{11} Among the following the correct statement(s) is(are)
\begin{qoptions}
\opt{A}{The principle involved in `zone refining' is fractional distillation}
\opt{B}{Parke's process in the metallurgy of silver is an application of Distribution Law}
\opt{C}{Preferential adsorption is the principle of froth-floatation process.}
\opt{D}{Open-Hearth furnace works on regenerative principle of heat economy.}
\end{qoptions}
\qn{12} Identify the incorrect statement(s)
\begin{qoptions}
\opt{A}{In acidic medium, both manganate and chromate ions undergo disproportionation.}
\opt{B}{Manganate ion is paramagnetic while permanganate ion is diamagnetic}
\opt{C}{Out of trivalent ions of Mn, Cr, V and Ti, the relatively more stable one is Mn$^{3+}$ ion}
\opt{D}{Higher oxides like Mn$_2$O$_7$, CrO$_3$ and V$_2$O$_5$ are predominantly basic in nature.}
\end{qoptions}
\qn{13} Which of the following statements is(are) correct ?
\begin{qoptions}
\opt{A}{XeO$_3$ + 2XeF$_6$ $\rightarrow$ 3 `x', the structure of x is square pyramidal}
\opt{B}{XeO$_3$ + XeOF$_4$ $\rightarrow$ 2 `y', the structure of y is see-saw}
\opt{C}{XeF$_2$, XeF$_4$, and XeF$_6$ react with H$_2$ to give Xe}
\opt{D}{XeF$_6$ does not react with SiO$_2$}
\end{qoptions}
\qn{14} Under certain conditions of concentration and temperature HNO$_3$ reacts with Zn and its reduction products are NO$_2$ and NO in a molar ratio 1 : 3. How many moles of HNO$_3$ are consumed by 1 mol of Zn?
\begin{qoptions}
\opt{A}{2.2}
\opt{B}{2.4}
\opt{C}{2.8}
\opt{D}{3.0}
\end{qoptions}
\sect{SECTION-I (iii)}
\paraq{15}{16}
The Diels-Alder reaction, a concerted [4+2]-cycloaddition between a diene and an olefin to yield a cyclohexene, was discovered in 1928 by Prof. Otto Diels and his coworker Kurt Alder. In Diels-Alder reaction the diene and the dienophile are functional groups in separate molecules. The Diels-Alder reaction is a concerted, one-step reaction that proceeds with high stereospecificity. The reaction can be intermolecular or intramolecular. An intramolecular generates two new rings in a single operation.\par
\qfig{p2_chem_s3_para15_16_1}
Diels-Alder reactions are reversible. Cyclohexenes can dissociate into a diene and an alkene by retro Diels-Alder reaction.
\qn{15} What is the structure of the intramolecular Diels-Alder product of the following compound?\par
\qfig{p2_chem_s3_q15_1}
\begin{qoptions}
\opt{A}{\qfigi{p2_chem_s3_q15_A}}
\opt{B}{\qfigi{p2_chem_s3_q15_B}}
\opt{C}{\qfigi{p2_chem_s3_q15_C}}
\opt{D}{\qfigi{p2_chem_s3_q15_D}}
\end{qoptions}
\qn{16} What compound would you use to prepare the following compound by an intramolecular Diels-Alder reaction?\par
\qfig{p2_chem_s3_q16_1}
\begin{qoptions}
\opt{A}{\qfigi{p2_chem_s3_q16_A}}
\opt{B}{\qfigi{p2_chem_s3_q16_B}}
\opt{C}{\qfigi{p2_chem_s3_q16_C}}
\opt{D}{\qfigi{p2_chem_s3_q16_D}}
\end{qoptions}
\paraq{17}{18}
Solvents have been classified in a number of ways depending on their physical and chemical properties. An inorganic nonaqueous solvent is a solvent other than water. Common examples are liquid ammonia, liquid sulphurdioxide, sulphuryl chloride, dinitrogen tetroxide antimony trichloride, bromine trifluoride, hydrogen fluoride etc. These are found to be electrically conducting and therefore expected to undergo self $-$ ionisation. These solvents are used in chemical research and industry for reactions that cannot occur in aqueous solutions.
\qn{17} Pick out correct statements among the following\par
(I) Dissolution of BrSbF$_8$ or KBrF$_4$ in liquid BrF$_3$ gives a highly conducting solution.\par
(II) BrSbF$_8$ acts as an acid in BrF$_3$\par
(III) KBrF$_4$ acts as a base in BrF$_3$\par
(IV) The reaction, AgBrF$_4$ + BrF$_2$SbF$_6$ $\rightarrow$ AgSbF$_6$ + 2BrF$_3$, is a neutralization reaction in BrF$_3$
\begin{qoptions}
\opt{A}{Only I and II}
\opt{B}{Only I and III}
\opt{C}{Only I ,II and III}
\opt{D}{All are correct}
\end{qoptions}
\qn{18} Which of the following statement is correct ?
\begin{qoptions}
\opt{A}{BrF$_5$ is strong fluorinating agent than BrF$_3$}
\opt{B}{KF acts as acid and SbF$_5$ acts as base in liquid BrF$_3$.}
\opt{C}{The liquid I$_2$Cl$_6$ has appreciable conductivity due to self ionisation into I$^{+}$, ICl$_6^{-}$}
\opt{D}{BrF$_5$ under goes disproportionation when dissolved in water}
\end{qoptions}
\raggedcolumns
\partsec{3}{MATHEMATICS}
\sect{SECTION-I (i)}
\qn{1} If z is a non real complex number satisfying $z^{26}=1$, then $\displaystyle\sum_{r=1}^{25}\dfrac{1}{1+z^{9r}+z^{18r}}=$
\begin{qoptions}
\opt{A}{14}
\opt{B}{15}
\opt{C}{16}
\opt{D}{17}
\end{qoptions}
\qn{2} Let the normal at $(2\cos\alpha,\sin\alpha)$ to the ellipse $x^2+4y^2=4$ intersects the ellipse again at $(2\cos\beta,\sin\beta)$ ($\alpha,\beta\in R$). If $\sin\beta=-\sin\alpha\left(\dfrac{a+b\sin^2\alpha}{c+d\sin^2\alpha}\right)$ in simplest form, then the value of $(a+b-c-d)$ is equal to ($a,b,c,d\in I$)
\begin{qoptions}
\opt{A}{$-3$}
\opt{B}{$-2$}
\opt{C}{0}
\opt{D}{4}
\end{qoptions}
\qn{3} If $I=\displaystyle\int_0^{4\pi}\log_e^{|13\sin x+3\sqrt3\cos x|}dx$, then the value of I is equal to
\begin{qoptions}
\opt{A}{$\log_e^7$}
\opt{B}{$2\pi\log_e^7$}
\opt{C}{$4\pi\log_e^7$}
\opt{D}{$8\pi\log_e^7$}
\end{qoptions}
\qn{4} Let $P_1(x)=\alpha x^2-\beta x-\gamma$, $P_2(x)=\beta x^2-\gamma x-\alpha$, $P_3(x)=\gamma x^2-\alpha x-\beta$ be three quadratic polynomials where $\alpha,\beta,\gamma$ are non-zero real numbers suppose there exists a real number k such that $P_1(k)=P_2(k)=P_3(k)$ then $\dfrac{9\alpha+8\beta+5\gamma}{5\alpha-3\beta-\gamma}=$\_\_\_\_\_\_
\begin{qoptions}
\opt{A}{19}
\opt{B}{15}
\opt{C}{21}
\opt{D}{22}
\end{qoptions}
\qn{5} Area enclosed by y = g(x), x = $-3$, x = 5 and x-axis where g(x) is the inverse function of $f(x)=x^3+3x+1$ is \_\_\_\_\_
\begin{qoptions}
\opt{A}{$\dfrac52$}
\opt{B}{$\dfrac72$}
\opt{C}{$\dfrac92$}
\opt{D}{$\dfrac32$}
\end{qoptions}
\qn{6} If $(\alpha_r,\beta_r)$ where = 1, 2, 3....... are the integral solutions of the equation $2x^2y^2+y^2-6x^2-12=0$, then the area of the polygon whose vertices are $(\alpha_r,\beta_r)$, r = 1, 2, 3.........is
\begin{qoptions}
\opt{A}{12 sq. units}
\opt{B}{16 sq. units}
\opt{C}{20 sq. units}
\opt{D}{18 sq. units}
\end{qoptions}
\sect{SECTION-I (ii)}
\qn{7} If a,b,c are positive real numbers such that a + 2b + 3c = 15 and the maximum value of 6(1+a)bc + a(2b+ 3c) is k, then which of the following is/are true ?
\begin{qoptions}
\opt{A}{The sum of the digits of k is 2}
\opt{B}{The total number of positive integral divisors of k is 12}
\opt{C}{The number formed by last two digits of k is divisible by 23}
\opt{D}{The total number of odd positive integral divisors of k is 3}
\end{qoptions}
\qn{8} If $I_n=\displaystyle\int_0^{\pi/2}x(\sin x+\cos x)^n dx$, and $\dfrac{nI_n-\frac\pi2}{I_{n-2}}=f(n)$, then which of the following is/are true ?
\begin{qoptions}
\opt{A}{The sum of the digits of f (2016) is 7}
\opt{B}{f (2017) is divisible by 63}
\opt{C}{f (2017) is divisible by 32}
\opt{D}{f (2016) is divisible by 31}
\end{qoptions}
\qn{9} If $f:\left(0,\dfrac\pi2\right)\to R^{+}$ be a real valued function satisfies the differential equation $xf'(x)=1+f(x)\left(x^2f(x)-1\right)$ and $f\left(\dfrac\pi4\right)=\dfrac4\pi$ and also $A=\displaystyle\int_0^{\pi/2}(\cos x)f(x)\,dx$ then which of the following is/are true ?
\begin{qoptions}
\opt{A}{$\dfrac2\pi<A<1$}
\opt{B}{$1<A<\dfrac\pi2$}
\opt{C}{$\dfrac76<A<\dfrac{\pi+6}{6}$}
\opt{D}{$\dfrac{1+\sqrt2}{2}<A<\dfrac{\pi+2\sqrt2}{4}$}
\end{qoptions}
\qn{10} If $P=\displaystyle\lim_{n\to\infty}\prod_{r=1}^{n}\left(1+\dfrac1{4r^2-1}\right)$, then [37P] is greater than or equal to, ([.] denotes greatest integer function)
\begin{qoptions}
\opt{A}{55}
\opt{B}{57}
\opt{C}{58}
\opt{D}{59}
\end{qoptions}
\qn{11} If P is a matrix of order 3X2 and Q is a matrix of order 2X3 such that $PQ=\begin{bmatrix}8&2&-2\\2&5&4\\-2&4&5\end{bmatrix}$ and also QP is an invertible matrix, then $|(QP)|\geq$ \_\_\_\_\_\_
\begin{qoptions}
\opt{A}{78}
\opt{B}{88}
\opt{C}{80}
\opt{D}{90}
\end{qoptions}
\qn{12} Let $f(x)=x^3+5x+8$ and $x=\alpha$ be a point such that $f'(\alpha)\neq\dfrac{f(b)-f(a)}{b-a}$ for any two values of a, b $\in$ R, then the number of such points is n, where n is less than
\begin{qoptions}
\opt{A}{1}
\opt{B}{2}
\opt{C}{3}
\opt{D}{4}
\end{qoptions}
\qn{13} If $^{29}C_1+5.^{29}C_5+9.^{29}C_9+\dots+29.^{29}C_{29}=k$, then k is divisible by
\begin{qoptions}
\opt{A}{27}
\opt{B}{23}
\opt{C}{32}
\opt{D}{64}
\end{qoptions}
\qn{14} If PA and PB are two tangents drawn from P to the parabola $4y^2=x$ and the normals drawn at A and B on the parabola are again intersect at a point Q on the same parabola. Then the locus of the centre of the circle passing through P, A and B is
\begin{qoptions}
\opt{A}{An Ellipse whose eccentricity is $\dfrac35$}
\opt{B}{A parabola whose length of latusrectums is $\dfrac1{32}$}
\opt{C}{An ellipse whose centre is $\left(\dfrac12,\dfrac32\right)$}
\opt{D}{A Parabola whose vertex is $\left(\dfrac14,0\right)$}
\end{qoptions}
\sect{SECTION-I (iii)}
\paraq{15}{16}
Let $f(x)=|||x|-1|-1|$ and $g(x)=\{\max(f(t);\ x\le t\le x+1,\ x\in R)\}$ then answer the following
\qn{15} The number of points of non-differentiability of g(x) is \_\_\_
\begin{qoptions}
\opt{A}{7}
\opt{B}{8}
\opt{C}{9}
\opt{D}{5}
\end{qoptions}
\qn{16} The graph of g (x) is symmetric about the line
\begin{qoptions}
\opt{A}{$x=\dfrac12$}
\opt{B}{$x=-\dfrac12$}
\opt{C}{x = 0}
\opt{D}{x = $-$ 1}
\end{qoptions}
\paraq{17}{18}
If a, b, c $>$ 0 are such that a + b + c = 3 and two positive numbers m, n $\neq$ 1 are such that $\max(a^m+b^m+c^m)=3$ and $\min(a^n+b^n+c^n)=3$ and a function $f(x)=x^2+px+q$ is such that f(x) = 0 has m, n as its roots, then answer the following
\qn{17} Which is correct ?
\begin{qoptions}
\opt{A}{(p + q +1)p $<$ 0}
\opt{B}{(p + q + 1)p $>$ 0}
\opt{C}{(p + q + 1)p = 0}
\opt{D}{Data Insufficient}
\end{qoptions}
\qn{18} If a, b, c are not all equal and $\alpha>0$ such that $a^\alpha+b^\alpha+c^\alpha=3$ and $f'(\alpha)=0$ then
\begin{qoptions}
\opt{A}{n = 2 $-$ m}
\opt{B}{n = 1 $-$ m}
\opt{C}{n = 3 $-$ m}
\opt{D}{m = n}
\end{qoptions}
\end{multicols}
\end{document}
